% *********************************************************************
% © 2010–2022 Jeremy Sylvestre
%
% Permission is granted to copy, distribute and/or modify this document
% under the terms of the GNU Free Documentation License, Version 1.3 or
% any later version published by the Free Software Foundation; with no
% Invariant Sections, no Front-Cover Texts, and no Back-Cover Texts. A
% copy of the license is included in the appendix entitled “GNU Free
% Documentation License” that appears in the output document of this
% PreTeXt source code. All trademarks™ are the registered® marks of
% their respective owners.
%
% *********************************************************************

\begin{tikzpicture}[
	vertex/.style={ circle, draw, very thin, fill, inner sep=0pt, minimum size=4pt },
	vertexg/.style={ black!40!white, circle, draw, very thin, fill, inner sep=0pt, minimum size=4pt },
	vertexo/.style={ circle, draw, very thin, inner sep=0pt, minimum size=4pt }
]

% very light "frame" so that his graph has roughly same width as
% components-overlap-2.tex
\draw[very thin,black!1!white] (-4.45,0) to (-4.5,0) to (-4.5,0.05);
\draw[very thin,black!1!white] (3.45,0) to (3.5,0) to (3.5,0.05);
\draw[very thin,black!1!white] (3.45,3.5) to (3.5,3.5) to (3.5,3.45);
\draw[very thin,black!1!white] (-4.45,3.5) to (-4.5,3.5) to (-4.5,3.45);

\node[vertex] at (-2,0) (g3v1) {};
\node[vertex] at (-1,1) (g3v2) {};
\draw (g3v2) to (g3v1);
\node[vertex] at (0,2) (g3v3) {};
\node[vertex] at (1,3) (g3v4) {};
\draw (g3v4) to (g3v3);
\draw (g3v2.north) to [out=90,in=180] (g3v3.west);
\draw (g3v2.east) to [out=0,in=270] (g3v3.south);
\draw (g3v4) arc (175:-175:0.25cm) to (g3v4);
\draw (g3v1) arc (5:355:0.25cm) to (g3v1);
\node[vertex] at (-2,3) (g3v5) {};
\node[vertex] at (-1,2) (g3v6) {};
\draw (g3v6) to (g3v5);
\node[vertex] at (0,1) (g3v7) {};

\node[vertex] at (1,0) (g3v8) {};
\draw (g3v8) to (g3v7);
\draw (g3v6.south) to [out=270,in=180] (g3v7.west);
\draw (g3v6.east) to [out=0,in=90] (g3v7.north);

\end{tikzpicture}
